\cvsection{Experience}

\begin{cventries}

  \cventry
    {Software Engineer $\to$ Senior Software Engineer, HyperIQ}
    {Cloudian Inc.}
    {Remote}
    {Mar. 2020 -- Jun. 2026}
    {
      \begin{cvitems}
        \item {Owned engineering for HyperIQ, a Grafana/Prometheus-based monitoring and observability platform shipped across 15 product releases and deployed on-premises in 400+ active customer environments alongside Cloudian HyperStore (S3-compatible object storage).}
        \item {Designed and shipped HyperIQ's alerting system (30+ production rules covering service availability, disk fill, SMART drive health, NTP drift, CPU temperature, network degradation, and Cassandra tombstone accumulation); later led full migration from Prometheus Alertmanager to Grafana unified alerting, preserving customer-customized thresholds and notification routing through upgrades spanning three major version gaps.}
        \item {Built \texttt{admin-exporter}, a custom Prometheus exporter in Go that queried the HyperStore Admin API to surface per-user, per-group, and per-bucket storage capacity metrics - complementing log-based request activity metrics and enabling quota visibility that logs alone couldn't provide; fixed a correctness bug where every EC-policy bucket reported the wrong raw size due to an overloaded API field, adding a full unit test matrix covering all five RF/EC storage policy configurations.}
        \item {Resolved S3 Analytics dashboards returning no data on high-cardinality clusters for time ranges over 7 days by replacing fixed-interval \texttt{rate()}/\texttt{increase()} windows with Grafana's \texttt{\$\_\_rate\_interval}, enabling enterprise deployments to visualize 30-day trends for the first time.}
        \item {Built multi-cluster monitoring support allowing a single HyperIQ instance to monitor multiple HyperStore clusters — extended \texttt{iqmgr} with cluster lifecycle commands (add/remove), propagated a unique \texttt{cluster} label across all Prometheus scrape configs, updated every Grafana dashboard query to filter by it, and persisted per-cluster config and servicemaps in etcd; preserved backward compatibility so existing timeseries without the label remained queryable after upgrade.}
        \item {Built \texttt{iqmgr config}, a configuration management CLI (get/set/list/apply) backed by etcd and confd templates, eliminating the problem of customer alert thresholds and notification settings (SMTP, Slack routing) being silently overwritten on every HIQ upgrade; configs persist in etcd and are applied to HyperIQ services via template rendering at apply time.}
        \item {Diagnosed and resolved P1 production incidents at enterprise accounts remotely from customer-provided logs and Prometheus data, shipping targeted fixes - sometimes as pre-release RPMs to unblock priority customers without waiting for the next release; also served as the primary technical authority for HyperIQ across SA and SE teams for pre-sales assessments and customer customization requests.}
        \item {Migrated the HyperIQ OVA base OS from Ubuntu to Rocky Linux 8.10, updating the full build pipeline, switching network configuration to NetworkManager, handling VMware/vSphere/Xen VM metadata requirements, and building installer logic for the cross-OS customer upgrade path.}
        \item {Built a Grafana Agent $\to$ OTel Collector config translation tool in Go for a customer deployment, generating valid receiver/exporter/pipeline configuration from existing scrape configs with TLS and bcrypt-based auth; packaged as an RPM with a CLI subcommand to update configs in place.}
      \end{cvitems}
    }

  \cventry
    {Software Engineer, HyperStore}
    {Cloudian Inc.}
    {San Mateo, CA}
    {Aug. 2017 -- Mar. 2020}
    {
      \begin{cvitems}
        \item {Maintained and extended \texttt{system\_setup.sh}, HyperStore's interactive node pre-configuration script, fixing reliability issues across disk partitioning/formatting, networking setup (including multi-NIC and VLAN environments), hostname configuration, SSH key management, and staging directory handling.}
        \item {Maintained HyperStore's Bash/Puppet/SaltStack installer, resolving reliability issues across node addition, multi-DC upgrade paths, and uninstall cleanup; improved upgrade reliability through survey file handling, conflicting servicemap detection, and correct settings migration across versions; automated TLS cert and keystore generation; extended HSH mode support with root password disable and SALT-based user management.}
      \end{cvitems}
    }

\end{cventries}
